\chapter{Background and Theoretical Foundations}
\label{ch:background}

%% ── 2.1 ─────────────────────────────────────────────────────────────────────
\section{LoRaWAN Technology}
\label{sec:lorawan-tech}

\subsection{LoRa Physical Layer}
% Chirp Spread Spectrum (CSS) modulation
% Spreading factors SF7–SF12: trade-off between range and data rate
% Frequency bands: EU868 (863–870 MHz), channels, bandwidth (125/250 kHz)

\subsection{LoRaWAN Protocol Architecture}
% Four-layer architecture: end device → gateway → network server → application server
% Class A operation (used by MKR WAN 1310): uplink-initiated, two Rx windows
% OTAA vs ABP activation — our project uses OTAA

\subsection{Duty Cycle and Payload Constraints}
% EU868: 1% duty cycle per sub-band
% Maximum payload size depends on SF and DR:
%   DR0 (SF12): 51 bytes — this is our constraint for FL payloads
%   DR5 (SF7): 222 bytes
% Implication: FL model updates MUST fit in 51 bytes at worst case
% This directly shaped our 52-byte payload design (4 header + 48 weights)

\subsection{The Things Network}
% TTN v3 infrastructure used in both prior and current project
% Webhook-based integration: TTN → fl_server.py
% Downlink limitations: ~10 confirmed downlinks/day per device

%% ── 2.2 ─────────────────────────────────────────────────────────────────────
\section{Indoor Radio Propagation and Path Loss}
\label{sec:path-loss-theory}

\subsection{Free-Space Path Loss}
% Friis equation: PL_fs = 20·log10(d) + 20·log10(f) + 20·log10(4π/c)
% Baseline for understanding signal attenuation

\subsection{Log-Distance Path Loss Model}
% PL(d) = PL(d_0) + 10·n·log10(d/d_0) + X_σ
% Path loss exponent n depends on environment (indoor office: n ≈ 2.0–3.5)
% Shadow fading X_σ: captures random variations

\subsection{Expected Path Loss Computation}
% How exp_pl is calculated in our dataset:
% exp_pl = TX_power - RSSI + antenna_corrections
% This is the target variable for our regression model

\subsection{Environmental Factors Affecting Indoor Propagation}
% Key insight from the prior project:
% - Humidity increases RF absorption (water molecule resonance)
% - Temperature gradients cause refraction
% - CO2/PM2.5 correlation with occupancy → obstacle effects
% - Pressure variations affect air density → propagation speed
% Reference to prior project findings and correlation analysis

%% ── 2.3 ─────────────────────────────────────────────────────────────────────
\section{Machine Learning for Path Loss Prediction}
\label{sec:ml-pathloss}

\subsection{Traditional Centralized Approaches}
% Overview: MLR, Random Forest, Gradient Boosting (XGBoost)
% Prior project results: XGBoost achieved R² ≈ 0.93
% Advantage: high accuracy with full centralized data
% Limitation: requires all data at one location

\subsection{Neural Network Approaches}
% Dense neural networks for regression
% Knowledge distillation: compressing ensemble models into small NNs
% Motivation for our Dense(5→8→1) architecture

%% ── 2.4 ─────────────────────────────────────────────────────────────────────
\section{Federated Learning}
\label{sec:fl-theory}

\subsection{Fundamentals}
% McMahan et al. (2017): "Communication-Efficient Learning of Deep Networks
%   from Decentralized Data"
% Core idea: train locally, aggregate globally, never share raw data
% Advantages: privacy, reduced communication, local adaptation

\subsection{Federated Averaging Algorithm}
\label{subsec:fedavg}

\begin{algorithm}[H]
  \caption{\gls{fedavg} --- Federated Averaging~\cite{mcmahan2017communication}}
  \label{alg:fedavg}
  \KwIn{$K$ clients, initial model $w_0$, $T$ rounds, $E$ local epochs, learning rate $\eta$}
  \KwOut{Trained global model $w_T$}
  \For{each round $t = 1, 2, \ldots, T$}{
    Server selects participating subset $S_t \subseteq \{1, \ldots, K\}$\;
    Server broadcasts current model $w_t$ to clients in $S_t$\;
    \For{each client $k \in S_t$ \textbf{in parallel}}{
      Initialize local model: $w_k \leftarrow w_t$\;
      \For{epoch $e = 1, \ldots, E$}{
        \For{each mini-batch $\mathcal{B} \subset \mathcal{D}_k$}{
          $w_k \leftarrow w_k - \eta \nabla \ell(w_k; \mathcal{B})$\;
        }
      }
      Client sends $w_k$ to server\;
    }
    Server aggregates: $\displaystyle w_{t+1} \leftarrow \sum_{k \in S_t}
      \frac{n_k}{\sum_{j \in S_t} n_j} \cdot w_k$\;
  }
\end{algorithm}

\subsection{Non-IID Data Challenges}
% In real deployments, each device sees different data distributions
% Our case: ED0 (stairwell) has different conditions than ED3 (kitchen)
% Non-IID causes slower convergence, potential divergence
% Strategies: larger local batches, more communication rounds

\subsection{Communication Efficiency Techniques}
% Quantization: float32 → int8 (our approach: scale factor 127)
% Gradient compression, sparse updates
% Critical for LoRaWAN: 51-byte payload limit at SF12

%% ── 2.5 ─────────────────────────────────────────────────────────────────────
\section{TinyML: Machine Learning on Microcontrollers}
\label{sec:tinyml-theory}

\subsection{TensorFlow Lite Micro}
% Designed for Cortex-M class processors
% Static memory allocation via tensor arena
% AllOpsResolver vs MutableOpResolver
% No dynamic memory, no filesystem, no OS dependency

\subsection{Model Quantization for Microcontrollers}
% Post-training quantization: float32 → int8
% Representative dataset for calibration
% Trade-off: accuracy vs model size vs inference speed

\subsection{On-Device Training Challenges}
% TFLite Micro does NOT support backpropagation natively
% Limited RAM prohibits storing full gradient computation graphs
% Practical approaches: simplified weight update rules, federated delta computation
% State of the art: research ongoing (TinyTrain, TinyOL)
